\newpage
\pagestyle{plain}

{\selectlanguage{russian}
\begin{center}
    \Large
    \textbf{Аңдатпа}
\end{center}

\vspace{0.5cm}

Фотоэлектрохимиялық (PEC) судың ыдырауы күн сәулесінен сутегі өндірудің ең перспективалы
жолдарының бірі болып табылады. Алайда жоғары өнімді фотоэлектрод материалдарын
эксперименттік іздеу баяу және ресурсты қажет ететін процесс болып қалуда.

Осы жұмыста материал мен эксперименттік параметрлер бойынша PEC ұяшықтарының үш
негізгі өнімділік көрсеткішін — фототок тығыздығын (mA/cm\textsuperscript{2}),
Күнге-Сутегіге (STH) тиімділігін (\%) және сутегінің бөліну жылдамдығын
($\mu$mol\,h\textsuperscript{-1}\,cm\textsuperscript{-2}) — болжайтын
машиналық оқыту жүйесі ұсынылған.

529 жарияланған эксперименттік нәтижеден жасалған 676 жазбадан тұратын деректер жиыны
физикалық бағалаулар мен синтетикалық жолдармен толықтырылды.
Жиырма физикаға негізделген белгі инженерлік тәсілмен жасалды.
XGBoost және Random Forest регрессорларының ансамблі мақсатты айнымалы бойынша
дербес оқытылды.

Ең жақсы кросс-валидация R\textsuperscript{2} көрсеткіштері:
Фототок үшін~— 0.454,
STH үшін~— 0.533,
H\textsubscript{2} үшін~— 0.366.

Жүйе Flask REST API және интерактивті веб-интерфейс арқылы ұсынылған.
}
